\section{September 4, 2026}

\subsection{Matrix--vector multiplication}

Let $V$ and $W$ be vector spaces over $\bb{F}$ with ordered bases
\[
  \mathcal B=(v_1,\ldots,v_n)
  \qquad\text{and}\qquad
  \mathcal C=(w_1,\ldots,w_m),
\]
respectively, and let $T\colon V\to W$ be linear. Recall that the matrix
representing $T$ with respect to these bases is
\[
  A=
  \begin{bmatrix}
    A_{11} & \cdots & A_{1n} \\
    \vdots &        & \vdots \\
    A_{m1} & \cdots & A_{mn}
  \end{bmatrix},
\]
where
\[
  T(v_j)=A_{1j}w_1+\cdots+A_{mj}w_m.
\]
Thus the $j$th column of $A$ is the coordinate vector
$[T(v_j)]_{\mathcal C}$.

Suppose that
\[
  v=x_1v_1+\cdots+x_nv_n,
  \qquad
  [v]_{\mathcal B}
  =\begin{bmatrix}x_1\\ \vdots\\ x_n\end{bmatrix}.
\]
By linearity,
\begin{align*}
  [T(v)]_{\mathcal C}
  &=x_1[T(v_1)]_{\mathcal C}+\cdots
    +x_n[T(v_n)]_{\mathcal C} \\
  &=x_1
    \begin{bmatrix}A_{11}\\ A_{21}\\ \vdots\\ A_{m1}\end{bmatrix}
    +\cdots+x_n
    \begin{bmatrix}A_{1n}\\ A_{2n}\\ \vdots\\ A_{mn}\end{bmatrix} \\
  &=
    \begin{bmatrix}
      A_{11}x_1+\cdots+A_{1n}x_n \\
      \vdots \\
      A_{m1}x_1+\cdots+A_{mn}x_n
    \end{bmatrix}.
\end{align*}
This motivates the following definition.

\begin{definition}[Matrix--vector product]
  Let $A=[A_{ij}]\in M_{m\times n}(\bb{F})$, and write its columns as
  $A_1,\ldots,A_n$. For
  \[
    x=\begin{bmatrix}x_1\\ \vdots\\ x_n\end{bmatrix}\in\bb{F}^n,
  \]
  define
  \[
    Ax=x_1A_1+\cdots+x_nA_n.
  \]
  Equivalently, the $k$th coordinate of $Ax$ is
  \[
    (Ax)_k=\sum_{j=1}^n A_{kj}x_j,
    \qquad k=1,\ldots,m.
  \]
\end{definition}

With this notation, the coordinate formula for a linear transformation is
\[
  [T(v)]_{\mathcal C}=A[v]_{\mathcal B}.
\]

\begin{example}
  Matrix--vector multiplication can be computed either as a linear
  combination of the columns or one coordinate at a time. For instance,
  \[
    \begin{bmatrix}
      1&2&3\\
      3&2&1
    \end{bmatrix}
    \begin{bmatrix}1\\2\\3\end{bmatrix}
    =1\begin{bmatrix}1\\3\end{bmatrix}
      +2\begin{bmatrix}2\\2\end{bmatrix}
      +3\begin{bmatrix}3\\1\end{bmatrix}
    =\begin{bmatrix}14\\10\end{bmatrix},
  \]
  while
  \[
    \begin{bmatrix}
      1&2&3\\
      4&5&6
    \end{bmatrix}
    \begin{bmatrix}1\\2\\3\end{bmatrix}
    =\begin{bmatrix}
      1\cdot1+2\cdot2+3\cdot3\\
      4\cdot1+5\cdot2+6\cdot3
    \end{bmatrix}
    =\begin{bmatrix}14\\32\end{bmatrix}.
  \]
\end{example}

\subsection{Finding matrix representations}

\begin{example}
  Define $T\colon\bb{R}^2\to\bb{R}^3$ by
  \[
    T\begin{bmatrix}x\\y\end{bmatrix}
    =\begin{bmatrix}x+2y\\2x-5y\\7y\end{bmatrix}.
  \]
  With respect to the standard bases, the columns of the representing
  matrix are
  \[
    T(e_1)=\begin{bmatrix}1\\2\\0\end{bmatrix}
    \qquad\text{and}\qquad
    T(e_2)=\begin{bmatrix}2\\-5\\7\end{bmatrix}.
  \]
  Hence
  \[
    A=
    \begin{bmatrix}
      1&2\\
      2&-5\\
      0&7
    \end{bmatrix}.
  \]
\end{example}

\begin{example}[Differentiation]
  Let $n\geq1$, and let
  \[
    T\colon\mathcal P_n(\bb{R})\longrightarrow\mathcal P_n(\bb{R}),
    \qquad T(p)=p',
  \]
  and use the ordered basis
  \[
    \mathcal B=(1,t,t^2,\ldots,t^n)
  \]
  for both the domain and codomain. Since
  \[
    T(1)=0,
    \qquad
    T(t)=1,
    \qquad
    T(t^j)=jt^{j-1},
  \]
  the matrix representing $T$ is
  \[
    A=
    \begin{bmatrix}
      0&1&0&\cdots&0\\
      0&0&2&\ddots&\vdots\\
      \vdots&\ddots&\ddots&\ddots&0\\
      0&\cdots&0&0&n\\
      0&\cdots&\cdots&0&0
    \end{bmatrix}.
  \]
\end{example}

\subsection{The vector space of linear maps}

Let $V$ and $W$ be vector spaces over $\bb{F}$. If
$T\colon V\to W$ is linear and $\alpha\in\bb{F}$, define
$\alpha T\colon V\to W$ pointwise by
\[
  (\alpha T)(v)=\alpha T(v).
\]
This map is again linear. Indeed, for $u,v\in V$,
\begin{align*}
  (\alpha T)(u+v)
  &=\alpha T(u+v) \\
  &=\alpha\bigl(T(u)+T(v)\bigr) \\
  &=(\alpha T)(u)+(\alpha T)(v),
\end{align*}
and, for $\beta\in\bb{F}$,
\begin{align*}
  (\alpha T)(\beta u)
  &=\alpha T(\beta u)
    =\alpha\beta T(u)
    =\beta(\alpha T)(u).
\end{align*}

Similarly, if $T_1,T_2\colon V\to W$ are linear, define their sum by
\[
  (T_1+T_2)(v)=T_1(v)+T_2(v).
\]
The map $T_1+T_2$ is linear because
\begin{align*}
  (T_1+T_2)(u+v)
  &=T_1(u+v)+T_2(u+v) \\
  &=(T_1+T_2)(u)+(T_1+T_2)(v)
\end{align*}
and
\[
  (T_1+T_2)(\alpha u)
  =\alpha(T_1+T_2)(u).
\]

\begin{definition}[Space of linear maps]
  The collection of all linear transformations from $V$ to $W$ is
  denoted by
  \[
    \mathcal L(V,W)
    =\{T\colon V\to W:T\text{ is linear}\}.
  \]
\end{definition}

\begin{proposition}
  The set $\mathcal L(V,W)$, equipped with pointwise addition and scalar
  multiplication, is a vector space over $\bb{F}$.
\end{proposition}

\begin{proof}
  The preceding calculations show that the two operations are closed in
  $\mathcal L(V,W)$. All vector-space identities follow pointwise from
  the corresponding identities in $W$. The additive identity is the
  zero transformation $0\colon V\to W$, defined by $0(v)=0$, and the
  additive inverse of $T$ is the map $-T$, defined by
  $(-T)(v)=-T(v)$.
\end{proof}

\subsection{Composition of linear maps}

Let $V$, $W$, and $Z$ be vector spaces over $\bb{F}$. Suppose that
\[
  T_2\colon V\longrightarrow W
  \qquad\text{and}\qquad
  T_1\colon W\longrightarrow Z
\]
are linear. Their \emph{composition} is the map
\[
  T_1\circ T_2\colon V\longrightarrow Z,
  \qquad
  (T_1\circ T_2)(v)=T_1\bigl(T_2(v)\bigr).
\]

\begin{proposition}
  The composition $T_1\circ T_2$ is linear.
\end{proposition}

\begin{proof}
  For $u,v\in V$ and $\alpha\in\bb{F}$,
  \begin{align*}
    (T_1\circ T_2)(u+v)
    &=T_1\bigl(T_2(u)+T_2(v)\bigr) \\
    &=T_1(T_2(u))+T_1(T_2(v)),
  \end{align*}
  and
  \[
    (T_1\circ T_2)(\alpha u)
    =T_1\bigl(\alpha T_2(u)\bigr)
    =\alpha T_1(T_2(u)).
  \]
\end{proof}
